% GENERATED, edit the French unit or the translation table
% GENERATED from content/data/enrichment.json by mxbuild.enrichment; do not edit.
\mxunit{enrichment-waiting-time-notes}{Waiting for the first success}
\mxsubtitle{Upper high school · probability and sequences \textperiodcentered{} Lesson and method \textperiodcentered{} 08}

\begin{prerequisites}
Independent trials, geometric sequences and sums.
\end{prerequisites}

\begin{keypoint}[Ideas and method]
Repeat independent trials, each with success probability p, where $0 < p \leq  1$. T is the trial number of the first success, starting at 1.

The event $T = k$ requires $k-1$ failures followed by success: $P(T=k) = (1-p)^{k-1}p$. Likewise, $P(T>n) = (1-p)^{n}$ for $n \geq  0$.

The expectation is $E(T)=\frac{1}{p}$. It follows from the tail-sum formula $E(T)=\sum  P(T>n)$, for $n \geq  0$, and the geometric series. This is an average, not a guaranteed success time.
\end{keypoint}

\begin{example}
For a fair die with success defined as rolling 6, $p=\frac{1}{6}$. The mean waiting time is 6 rolls, but the probability of waiting more than 6 rolls is $(\frac{5}{6})^{6}$, about $\num{0,335}$.
\end{example}

\begin{activity}[Explore]
For $p=\frac{1}{4}$, compare $P(T>n)$ for $n=0,1,\ldots ,10$. Relate the decrease to the geometric factor $\frac{3}{4}$.
\end{activity}

